%%%%%%%%%%First Chapter%%%%%%%%%%
\chapter{Introduction}

\section{Background}
Wireless communication systems continue to require higher data rates, wider coverage, and more reliable link quality. These requirements become more challenging as communication systems operate at higher frequencies, where propagation loss, blockage, antenna alignment, and coverage limitation can strongly affect the received signal. In this condition, the antenna is not only a radiating component, but also an important part of the overall communication channel because its gain, radiation pattern, polarization, and beam direction directly influence the received power and link performance.

Directional antennas are widely used to improve link quality by concentrating electromagnetic energy toward a desired direction. However, conventional high-gain antenna structures can be bulky, complex, or difficult to integrate into compact communication platforms. A metasurface transmitarray lens antenna offers an attractive alternative because it can provide spatial phase control using a planar arrangement of engineered unit cells. By controlling the transmission phase across the aperture, the transmitarray lens can transform the incident wave from a feed antenna into a focused or collimated wavefront, which can improve gain and support directional communication.

The performance of a metasurface transmitarray lens antenna depends on both the electromagnetic behavior of its unit cells and the system-level effect of the resulting radiation pattern. Therefore, antenna-level evaluation alone is not sufficient to understand its usefulness in a communication scenario. Parameters such as transmission magnitude, phase response, gain, directivity, beam direction, and side-lobe level must be connected with channel-related metrics such as received power, path loss, signal-to-noise ratio, and coverage. This relationship is important because an antenna design that performs well in isolation may still produce different link behavior depending on propagation distance, antenna orientation, and environmental assumptions.

Based on this background, this thesis investigates a metasurface transmitarray lens antenna through simulation and communication channel evaluation. ANSYS HFSS is used to evaluate the electromagnetic characteristics of the antenna, while channel-level analysis is used to observe how the antenna performance affects the communication link. The study focuses on building a simulation workflow that connects the physical antenna design with practical communication performance indicators.

\section{Motivation}
The main motivation of this research is the need to evaluate antenna designs from both electromagnetic and communication perspectives. In many antenna studies, the evaluation is mainly focused on parameters such as reflection coefficient, gain, directivity, and radiation pattern. Although these parameters are essential, they do not fully describe how the antenna affects the received signal in an actual communication link. For directional antennas such as metasurface transmitarray lenses, this gap becomes important because beam focusing, phase compensation, and aperture efficiency can directly change the channel response.

Another motivation is the practical need for compact and efficient high-gain antenna structures. A metasurface transmitarray lens antenna can provide beam-forming capability with a relatively low-profile planar structure. This makes it a promising candidate for communication systems that require directional radiation while maintaining simpler integration compared with bulky lens or reflector structures. However, the design must be carefully evaluated because unit cell loss, limited phase coverage, feed illumination, and phase error can reduce the final antenna performance.

This thesis is also motivated by the need for a clear simulation procedure that links unit cell design, full antenna simulation, and communication channel evaluation. By combining these stages, the antenna can be assessed not only by its radiation characteristics but also by its influence on link-related metrics such as received power, path loss, signal-to-noise ratio, and coverage. This approach provides a more complete understanding of whether the designed metasurface transmitarray lens antenna can support the intended communication scenario.

\section{Thesis Structure}
This thesis is organized as follows. Chapter 1 introduces the research background, motivation, and thesis structure. Chapter 2 presents the theoretical background of metasurface transmitarray lens antennas and communication channel evaluation. Chapter 3 describes the simulation setup and design workflow, including unit cell design, transmitarray lens configuration, ANSYS HFSS simulation, and channel evaluation setup. Chapter 4 presents the simulation results and discussion for both antenna-level and communication-level performance. Chapter 5 summarizes the conclusions, main contributions, limitations, and future work.
